%______________________________________________________________________
%  To generate the pdf:
%    pdflatex LODI_reflectionCoeff_tests.tex
%    pdflatex LODI_reflectionCoeff_tests.tex   (run twice to resolve
%                                                cross-references/citations)
%______________________________________________________________________
\documentclass[11pt]{article}
\usepackage{amsmath}
\usepackage{amssymb}
\usepackage{graphicx}
\usepackage{booktabs}
\usepackage[margin=1in]{geometry}
\usepackage{hyperref}
\usepackage{xcolor}

\newcommand{\sus}{\texttt{sus}}
\newcommand{\ups}{\texttt{.ups}}

\title{LODI Outflow Boundary Condition: \\
       Reflection-Coefficient Validation Tests}
\author{Todd Harman}
\date{October 2026}

\begin{document}
\maketitle

%______________________________________________________________________
\section{Motivation}
\label{sec:motivation}

The LODI (Local One-Dimensional Inviscid) outflow boundary condition in
\texttt{CCA/\allowbreak Components/\allowbreak ICE/\allowbreak CustomBCs/\allowbreak LODI2.cc} relaxes the boundary pressure
toward a reference value $p_\infty$ using a user-specified coefficient
$\sigma$, following the NSCBC (Navier--Stokes Characteristic Boundary
Conditions) formulation of Sutherland and Kennedy~\cite{SutherlandKennedy2003},
itself building on Rudy and Strikwerda~\cite{RudyStrikwerda1980} and Poinsot
and Lele~\cite{PoinsotLele1992}. A larger $\sigma$ anchors the boundary
pressure more strongly (suppressing long-time drift away from $p_\infty$) at
the cost of reflecting more of any genuine outgoing disturbance back into the
domain -- the central trade-off in the NSCBC literature.

A review of \texttt{LODI2.cc} raised the question of whether a factor of
$0.5$ appearing twice in the code (once in the default characteristic-wave
definitions $L_1,L_5$, once again in the outflow relaxation term) constituted
a double-counted, spurious halving of $\sigma$'s effective strength. Reading
the cited reference directly (Sutherland and Kennedy, Table 5 and Eq.\
8--9) showed both factors of $0.5$ are part of that paper's own convention
and are not duplicated -- the code matches its cited source exactly. This
document describes the follow-up validation effort: a suite of test problems
that directly measures the boundary's \emph{reflection coefficient} $|R|$ as
a function of $\sigma$, giving an empirical, reproducible check of the
boundary's behavior that is independent of the symbolic algebra review.

%______________________________________________________________________
\section{Background: the outflow relaxation term}
\label{sec:background}

For a subsonic outflow with the boundary-normal direction aligned with a
coordinate axis, \texttt{LODI2.cc} sets the incoming characteristic wave
amplitude ($L_1$ at a right/$+$ boundary, or $L_5$ at a left/$-$ boundary) to
\begin{equation}
  L_{\text{relax}} = \frac{1}{2L}\,\sigma\, c\,\bigl(1 - M^2\bigr)\,
                      \bigl(p - p_\infty\bigr),
  \label{eq:relax}
\end{equation}
where $c$ is the local speed of sound, $M$ is the local Mach number, $L$ is
the domain length normal to the boundary, and $p_\infty$ is the target
(ambient) pressure -- this is Sutherland and Kennedy's Eq.\ 8--9 verbatim.
$\sigma=0$ reduces to a purely characteristic (non-reflecting) extrapolation;
large $\sigma$ pulls the boundary pressure rigidly toward $p_\infty$, which
behaves increasingly like a reflecting wall. The tests below measure where a
given $\sigma$ actually sits on that spectrum.

%______________________________________________________________________
\section{Test methodology}
\label{sec:methodology}

%__________________________________
\subsection{Riemann-invariant decomposition at a probe}

Each test initializes a small-amplitude Gaussian temperature spike at rest,
off-center in the domain, using \texttt{customInitialization}
$\rightarrow$ \texttt{gaussianTemperature}. A disturbance introduced this way
decomposes into outward-traveling acoustic wave packets plus a stationary
entropy anomaly that never reaches a distant probe. At a probe location
on the domain's symmetry line (defined precisely for each test family
below), the probe's $\texttt{press\_CC}$ and the velocity component normal
to the boundary under test, $u$, are combined into right- and left-going
Riemann invariants
\begin{align}
  p_{+}(t) &= \tfrac{1}{2}\bigl(p(t)-p_\infty\bigr) + \tfrac{1}{2}\rho_0 c_0\, u(t)
              \qquad \text{(incident, right-going)}, \\
  p_{-}(t) &= \tfrac{1}{2}\bigl(p(t)-p_\infty\bigr) - \tfrac{1}{2}\rho_0 c_0\, u(t)
              \qquad \text{(reflected, left-going)},
  \label{eq:riemann}
\end{align}
where $\rho_0,c_0$ are the ambient density and sound speed. The incident wave
appears as a peak in $p_+(t)$ near the time the pulse first reaches the
probe; the boundary's reflection appears as a later peak in $p_-(t)$. The
reflection coefficient reported for each test is
\begin{equation}
  |R| \;=\; \frac{\max_t |p_-(t)|}{\max_t |p_+(t)|} .
  \label{eq:Rdef}
\end{equation}
This decomposition is exact only where the transverse velocity component is
identically zero; both the 1D and 2D test families below are constructed so
that the probe sits exactly on such a line.

%__________________________________
\subsection{On-the-fly probe extraction}

The probe history is recorded by the \texttt{DataAnalysis}
$\rightarrow$ \texttt{lineExtract} on-the-fly analysis module
(\texttt{CCA/\allowbreak Components/\allowbreak OnTheFlyAnalysis/\allowbreak lineExtract.cc}), sampling
\texttt{press\_CC} and \texttt{vel\_CC} at the probe cell every timestep into
one small, continuously-appended text file.

The post-processing itself is \texttt{\footnotesize parametricStudies/postProcessTools/compute\_LODI\_reflectionCoeff.py}, which reads the
\texttt{lineExtract} output directly (no further simulation queries needed),
forms Eqs.~\eqref{eq:riemann}--\eqref{eq:Rdef}, and writes $|R|$ to the
parametric-study framework's result file.

%__________________________________
\subsection{1D test family}
\label{sec:methodology-1D}

\texttt{\footnotesize StandAlone/inputs/ICE/LODI\_reflection/Lodi\_reflectionCoeff\_d\{x,y,z\}.ups}
each use a long, thin domain (1D in the named axis; the other two directions
are a single cell with symmetric boundaries) of length $L=2.0$~m. The pulse
sits at $x_{\text{spike}}=1.2$~m; the probe sits at $x_{\text{probe}}=1.6$~m
(and analogously for $y,z$). Because the pulse is centered off-axis, it
radiates two equal-amplitude planar wave packets, one toward each boundary.
The far boundary is also LODI, so with
\begin{align}
  t_1 &= \frac{x_{\text{probe}} - x_{\text{spike}}}{c_0}
         = 1.48\times10^{-3}\ \text{s} &&\text{(incident arrival)}, \\
  t_2 &= \frac{2L - x_{\text{spike}} - x_{\text{probe}}}{c_0}
         = 4.44\times10^{-3}\ \text{s} &&\text{(reflection being measured)}, \\
  t_3 &= \frac{x_{\text{spike}} + x_{\text{probe}}}{c_0}
         = 10.4\times10^{-3}\ \text{s} &&\text{(far-boundary contamination)},
  \label{eq:timing1D}
\end{align}
($c_0 = \sqrt{\gamma(\gamma-1)c_v T_0} = 270.12$~m/s for this test gas),
the ordering $t_1 < t_2 \ll t_3$ gives a clean, uncontaminated window around
$t_2$ in which to find the reflected peak.

%__________________________________
\subsection{2D test family}
\label{sec:methodology-2D}

\texttt{Lodi\_reflectionCoeff\_2D\_\{xy,xz,yz\}.ups} replace the 1D family's
planar pulse with a genuinely two-dimensional, radially-expanding
(cylindrical) one: both in-plane spreads of the Gaussian are finite, so the
wavefront that reaches the boundary under test is curved, striking it at a
continuum of incidence angles rather than head-on. This directly exercises
the transverse terms that the 1D LODI approximation in
Eq.~\eqref{eq:relax} drops.

The pulse and the probe both sit on the domain's centerline in the
\emph{second} active axis (e.g., for the x-y plane test, both are at
$y=L/2$); by the mirror symmetry of a centered pulse, the velocity
component along that second axis is exactly zero along the whole
centerline, which is what keeps the simple decomposition of
Eq.~\eqref{eq:riemann} exact even though the wave itself is curved. All four
in-plane boundaries are LODI; the domain size and pulse/probe placement
($L=1.0$~m, spike at $0.6$~m, probe at $0.8$~m, both on the $0.5$~m
centerline) are chosen, via the same image-source construction used for
Eq.~\eqref{eq:timing1D}, so that the two other boundaries' reflections
arrive safely after the one being measured:
\begin{equation}
  t_1 = 7.40\times10^{-4}\ \text{s} \;<\;
  t_2 = 2.22\times10^{-3}\ \text{s} \;<\;
  t_{\text{other}} = 3.78\times10^{-3}\ \text{s} \;<\;
  t_3 = 5.18\times10^{-3}\ \text{s}.
  \label{eq:timing2D}
\end{equation}

The x-z and y-z plane variants are the identical construction with the
in-plane axes relabeled; the x-z plane re-tests the x$+$ boundary (a
cross-check of the implementation under a different domain orientation and
patch decomposition), while the y-z plane tests the y$+$ boundary (a
genuinely different boundary face). By the test gas's isotropy, all three
should -- and, as reported below, do -- land on the same $|R|(\sigma)$
curve.

%______________________________________________________________________
\section{Results}
\label{sec:results}

%__________________________________
\subsection{1D (planar incidence)}

Table~\ref{tab:1D} gives $|R|(\sigma)$ for the 1D family. The x, y, and z
variants are bit-for-bit identical at every $\sigma$, as expected since they
are the same physics problem with the active axis permuted.

\begin{table}[h]
\centering
\caption{1D reflection coefficient vs.\ $\sigma$ (x, y, z variants agree
         exactly).}
\label{tab:1D}
\begin{tabular}{@{}cc@{}}
\toprule
$\sigma$ & $|R|$ \\
\midrule
0.00 & 0.01386 \\
0.05 & 0.01415 \\
0.10 & 0.01444 \\
0.27 (default) & 0.01546 \\
0.50 & 0.01700 \\
1.00 & 0.02116 \\
2.00 & 0.04060 \\
\bottomrule
\end{tabular}
\end{table}

%__________________________________
\subsection{2D (curved incidence)}

Table~\ref{tab:2D} gives $|R|(\sigma)$ for the 2D family. The x-y and y-z
plane sweeps are bit-for-bit identical at every $\sigma$ tested; the x-z
plane was spot-checked at $\sigma=0.27$ and also agrees exactly
($|R|=0.026309$). This is a nontrivial cross-check: the y-z plane result in
particular tests a genuinely different boundary face (y$+$ instead of x$+$)
in a different domain orientation, yet the test gas's isotropy still drives
all three to the same curve.

\begin{table}[h]
\centering
\caption{2D reflection coefficient vs.\ $\sigma$ (x-y and y-z plane sweeps
         agree exactly; x-z spot-checked at $\sigma=0.27$).}
\label{tab:2D}
\begin{tabular}{@{}cc@{}}
\toprule
$\sigma$ & $|R|$ \\
\midrule
0.00 & 0.02436 \\
0.05 & 0.02467 \\
0.10 & 0.02505 \\
0.27 (default) & 0.02631 \\
0.50 & 0.02809 \\
1.00 & 0.03255 \\
2.00 & 0.04477 \\
\bottomrule
\end{tabular}
\end{table}

%__________________________________
\subsection{Discussion}

Both families show the reflection coefficient increasing monotonically with
$\sigma$, consistent with the trade-off described in
Section~\ref{sec:background}: the default $\sigma=0.27$ keeps $|R|$ to a few
percent, while $\sigma=2$ pushes $|R|$ toward $4$--$5\%$ in both families.

The 2D family reflects systematically more than the 1D family at every
$\sigma$ -- roughly $70\%$ more at the default $\sigma=0.27$
($|R|=0.0263$ vs.\ $0.0155$). This is a real, physically meaningful result,
not an inconsistency: the 1D LODI approximation underlying
Eq.~\eqref{eq:relax} is derived assuming a locally one-dimensional,
normal-incidence wave, and drops the transverse pressure/velocity gradients
that a genuinely curved wavefront has. The curved-wave test exercises
exactly the terms the approximation omits, so a larger reflection
coefficient there is the expected signature of that approximation's known
limitation, not a defect in the implementation.

%______________________________________________________________________
\section{Conclusions}

\begin{itemize}
\item The LODI outflow boundary's $\sigma$ relaxation term behaves as the
      NSCBC literature predicts: reflection increases monotonically with
      $\sigma$, and the default $\sigma=0.27$ keeps reflection to a few
      percent for normal-incidence (1D) waves.
\item A genuinely curved (2D) wavefront reflects noticeably more than a
      planar one at every $\sigma$ tested -- an expected consequence of the
      1D LODI approximation, confirmed quantitatively rather than assumed.
\item Three independent 2D plane orientations, including two genuinely
      different boundary faces (x$+$ and y$+$), converge to the same
      $|R|(\sigma)$ curve, which is strong evidence against any
      axis-specific implementation bug.
\end{itemize}

\begin{thebibliography}{9}

\bibitem{SutherlandKennedy2003}
J.~C. Sutherland and C.~A. Kennedy,
``Improved boundary conditions for viscous, reacting, compressible flows,''
\emph{Journal of Computational Physics}, vol.~191, no.~2, pp.~502--524, 2003.
\url{https://doi.org/10.1016/j.jcp.2003.06.001}

\bibitem{PoinsotLele1992}
T.~J. Poinsot and S.~K. Lele,
``Boundary conditions for direct simulations of compressible viscous flows,''
\emph{Journal of Computational Physics}, vol.~101, no.~1, pp.~104--129, 1992.
\url{https://doi.org/10.1016/0021-9991(92)90046-2}

\bibitem{RudyStrikwerda1980}
D.~H. Rudy and J.~C. Strikwerda,
``A nonreflecting outflow boundary condition for subsonic Navier-Stokes
calculations,''
\emph{Journal of Computational Physics}, vol.~36, no.~1, pp.~55--70, 1980.
\url{https://doi.org/10.1016/0021-9991(80)90174-6}

\bibitem{YooIm2007}
C.~S. Yoo and H.~G. Im,
``Characteristic boundary conditions for simulations of compressible
reacting flows with multi-dimensional, viscous and reaction effects,''
\emph{Combustion Theory and Modelling}, vol.~11, no.~2, pp.~259--286, 2007.
\url{https://doi.org/10.1080/13647830600898995}

\end{thebibliography}

\end{document}
